% GENERATED FILE. Edit canonical lessons, then run npm run generate:books.
% canonical-source-hash: fnv1a64:a8126b16f0ac534d
% canonical-lessons: AR-C106-imkaniyya, AR-C106-sulta, AR-C106-najah, AR-C106-itqan, AR-C106-qudra

\chapter{Possibility, Authority, Success, Mastery, Ability}
\label{ch:ar-possibility-authority-success-mastery-ability}
\hlchaptermodality{106}

\begin{chapteropening}
\textbf{By the end of this chapter:} \emph{I can say a possibility, authority, success, mastery and ability.}

Complete the last lesson of chapter 106: I can say a possibility, authority, success, mastery and ability.
\end{chapteropening}

\section[\texorpdfstring{\ar{إمكانية} (imkāniyya)}{إمكانية (imkāniyya)}]{\ar{إمكانية} (imkāniyya) — a possibility}

\label{lesson:AR-C106-imkaniyya}



\begin{quote}
Before the new one: say the Arabic for around, then the Arabic for at.
\end{quote}

\subsection*{You'll want to know: \ar{إمكانية}}
\textbf{\ar{إمكانية}} — \emph{imkāniyya} — \textquotedblleft{}a possibility\textquotedblright{}.

\begin{cousinweb}
\textbf{\ar{إمكانية}} (\emph{imkāniyya}) is built on the native Arabic root \emph{m-k-n}, \textquotedblleft{}to have power, to be firm in place\textquotedblright{} (Lane). \emph{Imkāniyya} is a modern abstract noun from \emph{imkān}, \textquotedblleft{}possibility\textquotedblright{}, the verbal noun of Form IV \emph{amkana}, \textquotedblleft{}to be possible\textquotedblright{}.
\end{cousinweb}

\begin{grammarlens}[title={What you've built}]
One more word for this chapter.
\end{grammarlens}

\subsection*{Your turn}
\begin{itemize}
\raggedright
  \item \emph{Say it:} \emph{imkāniyya}
  \item \emph{Say it:} \emph{imkāniyya}, once more
  \item \emph{Say it:} say \emph{imkāniyya}
  \item \emph{Recall:} say the Arabic for around, then the Arabic for at, then say \emph{imkāniyya} again
\end{itemize}

\begin{tcolorbox}[breakable,colback=teal!4,colframe=teal!35!black,title={Before you move on}]
What does \ar{إمكانية} mean? (\textquotedblleft{}A possibility\textquotedblright{}.) Say it once more.
\end{tcolorbox}
\section[\texorpdfstring{\ar{سلطة} (sulṭa)}{سلطة (sulṭa)}]{\ar{سلطة} (sulṭa) — authority}

\label{lesson:AR-C106-sulta}



\begin{quote}
Before the new one: say the Arabic for at, then the Arabic for a possibility.

Answer \emph{masāʾ al-khayr}. (\emph{Masāʾ an-nūr}: the same shape as the morning reply.)
\end{quote}

\subsection*{You'll want to know: \ar{سلطة}}
\textbf{\ar{سلطة}} — \emph{sulṭa} — \textquotedblleft{}authority\textquotedblright{}.

\begin{cousinweb}
\textbf{\ar{سلطة}} (\emph{sulṭa}) is built on the native Arabic root \emph{s-l-ṭ}, \textquotedblleft{}to have mastery\textquotedblright{} (Lane). The same root gives \emph{sulṭān}, \textquotedblleft{}power, ruler\textquotedblright{}, the source of English \emph{sultan}.
\end{cousinweb}

\begin{grammarlens}[title={What you've built}]
One more word for this chapter.
\end{grammarlens}

\subsection*{Your turn}
\begin{itemize}
\raggedright
  \item \emph{Say it:} \emph{sulṭa}
  \item \emph{Say it:} \emph{sulṭa}, once more
  \item \emph{Say it:} say \emph{sulṭa}
  \item \emph{Recall:} say the Arabic for at, then the Arabic for a possibility, then say \emph{sulṭa} again
\end{itemize}

\begin{tcolorbox}[breakable,colback=teal!4,colframe=teal!35!black,title={Before you move on}]
What does \ar{سلطة} mean? (\textquotedblleft{}Authority\textquotedblright{}.) Say it once more.
\end{tcolorbox}
\section[\texorpdfstring{\ar{نجاح} (najāḥ)}{نجاح (najāḥ)}]{\ar{نجاح} (najāḥ) — success}

\label{lesson:AR-C106-najah}



\begin{quote}
Before the new one: say the Arabic for a possibility, then the Arabic for authority.
\end{quote}

\subsection*{You'll want to know: \ar{نجاح}}
\textbf{\ar{نجاح}} — \emph{najāḥ} — \textquotedblleft{}success\textquotedblright{}.

\begin{cousinweb}
\textbf{\ar{نجاح}} (\emph{najāḥ}) is built on the native Arabic root \emph{n-j-ḥ}, \textquotedblleft{}to succeed, to prosper\textquotedblright{} (Lane). \emph{Najāḥ} is the verbal noun, \textquotedblleft{}success\textquotedblright{}.
\end{cousinweb}

\begin{grammarlens}[title={What you've built}]
One more word for this chapter.
\end{grammarlens}

\subsection*{Your turn}
\begin{itemize}
\raggedright
  \item \emph{Say it:} \emph{najāḥ}
  \item \emph{Say it:} \emph{najāḥ}, once more
  \item \emph{Say it:} say \emph{najāḥ}
  \item \emph{Recall:} say the Arabic for a possibility, then the Arabic for authority, then say \emph{najāḥ} again
\end{itemize}

\begin{tcolorbox}[breakable,colback=teal!4,colframe=teal!35!black,title={Before you move on}]
What does \ar{نجاح} mean? (\textquotedblleft{}Success\textquotedblright{}.) Say it once more.
\end{tcolorbox}
\section[\texorpdfstring{\ar{إتقان} (itqān)}{إتقان (itqān)}]{\ar{إتقان} (itqān) — mastery}

\label{lesson:AR-C106-itqan}



\begin{quote}
Before the new one: say the Arabic for authority, then the Arabic for success.
\end{quote}

\subsection*{You'll want to know: \ar{إتقان}}
\textbf{\ar{إتقان}} — \emph{itqān} — \textquotedblleft{}mastery\textquotedblright{}.

\begin{cousinweb}
\textbf{\ar{إتقان}} (\emph{itqān}) is built on the native Arabic root \emph{t-q-n} (Lane). \emph{Itqān} is the verbal noun of Form IV \emph{atqana}, \textquotedblleft{}to make firm, to do skilfully\textquotedblright{}.
\end{cousinweb}

\begin{grammarlens}[title={What you've built}]
One more word for this chapter.
\end{grammarlens}

\subsection*{Your turn}
\begin{itemize}
\raggedright
  \item \emph{Say it:} \emph{itqān}
  \item \emph{Say it:} \emph{itqān}, once more
  \item \emph{Say it:} say \emph{itqān}
  \item \emph{Recall:} say the Arabic for authority, then the Arabic for success, then say \emph{itqān} again
\end{itemize}

\begin{tcolorbox}[breakable,colback=teal!4,colframe=teal!35!black,title={Before you move on}]
What does \ar{إتقان} mean? (\textquotedblleft{}Mastery\textquotedblright{}.) Say it once more.
\end{tcolorbox}
\section[\texorpdfstring{\ar{قدرة} (qudra)}{قدرة (qudra)}]{\ar{قدرة} (qudra) — ability}

\label{lesson:AR-C106-qudra}



\begin{quote}
Before the new one: say the Arabic for success, then the Arabic for mastery.
\end{quote}

\subsection*{You'll want to know: \ar{قدرة}}
\textbf{\ar{قدرة}} — \emph{qudra} — \textquotedblleft{}ability\textquotedblright{}.

\begin{cousinweb}
\textbf{\ar{قدرة}} (\emph{qudra}) is built on the native Arabic root \emph{q-d-r}, \textquotedblleft{}to measure; to have power\textquotedblright{} (Lane). \emph{Qudra} is power, the ability to do.
\end{cousinweb}

\begin{grammarlens}[title={What you've built}]
One more word for this chapter.
\end{grammarlens}

\subsection*{Your turn}
\begin{itemize}
\raggedright
  \item \emph{Say it:} \emph{qudra}
  \item \emph{Say it:} \emph{qudra}, once more
  \item \emph{Say it:} say \emph{qudra}
  \item \emph{Recall:} say the Arabic for success, then the Arabic for mastery, then say \emph{qudra} again
\end{itemize}

\begin{tcolorbox}[breakable,colback=teal!4,colframe=teal!35!black,title={Before you move on}]
What does \ar{قدرة} mean? (\textquotedblleft{}Ability\textquotedblright{}.) Say it once more.
\end{tcolorbox}
